\documentclass[11pt]{article}
\usepackage[english]{babel}
\usepackage[margin=0.82in]{geometry}
\usepackage{amsmath,amssymb,booktabs,tabularx,array}
\usepackage[hidelinks]{hyperref}

\setlength{\parindent}{0pt}
\setlength{\parskip}{0.55\baselineskip}
\setlength{\emergencystretch}{3em}
\newcolumntype{L}[1]{>{\raggedright\arraybackslash}p{#1}}

\title{Technical Review of the Public Emergency Codes\\
Mortality and Economic Manuscripts}
\author{}
\date{August 8, 2026}

\begin{document}
\maketitle

\section*{Overall assessment}

The two manuscripts are unusually explicit about the distinction between an
identified PEC effect and an assumption-conditional planning calculation.  The
headline mortality calculations, modal values, and Monte Carlo summaries were
rerun from the archived scripts and reproduced numerically.  The key quoted
values from the Nguyen bystander-CPR study, the randomized video-dispatch
trial, the 2026 HHS VSL/VQALY guidance, and the OMB A-94 reinstatement were also
spot-checked against the cited primary sources and were quoted accurately.

Major revision is nevertheless needed before the economic results can be
described as enforcing the paper's claimed accounting and incident-level
safeguards.  Several safeguards are present only as prose or as algebraic
identities that are true by construction; the executable model does not
implement them.  The mortality paper is more internally coherent, but its
annual-count implementation does not preserve the claimed cross-pathway
exclusivity, and its CPR time conversion should be labeled as a stronger
approximation than the current wording suggests.

\section*{Major technical findings}

\subsection*{1. Definite error: the fiscal/capacity ``partition'' is tautological in code}

\paragraph{Location.}
\texttt{Main\_economic.tex}.\\
``Economic perspectives and definitions,''
Equation~\eqref{eq:publicpartition-review} below, and
\path{pec_economic_simulation.py} in the per-pathway public-resource block.

The manuscript states
\begin{equation}
G_{ir}=F_{ir}+C^{\mathrm{res}}_{ir},\qquad
C^{\mathrm{res}}_{ir}=G_{ir}-F_{ir},
\tag{as printed in the manuscript}
\label{eq:publicpartition-review}
\end{equation}
and says that the code enforces a mutually exclusive allocation of a common
gross public resource.  The code instead samples public fiscal value and
capacity value independently, defines gross value as their sum, and then
subtracts fiscal value to recover the original capacity draw.  The subsequent
assertion verifies only the identity
\(F+C=(F+C)\); it does not test whether the two sampled components value the
same dispatcher hour, responder hour, transport, bed-hour, or medical resource
twice.

\paragraph{Why it matters.}
The mature public-resource headline of approximately \(\$1.173\) billion and
its split into \(\$943\) million fiscal plus \(\$230\) million capacity are
presented as non-overlapping.  The current implementation does not establish
that property, so the split and combined total may contain conceptual overlap.

\paragraph{Suggested fix.}
Model one gross resource quantity from operational primitives and allocate it
with an explicit share, for example
\[
F_{ir}=\alpha_{ir}G_{ir},\qquad
C^{\mathrm{res}}_{ir}=(1-\alpha_{ir})G_{ir},\qquad 0\leq\alpha_{ir}\leq1,
\]
or calculate fiscal and residual capacity from a common unit ledger.  Assert
the identity against that independently calculated gross quantity, then rerun
all public-system, pathway, feature, and headline tables.

\subsection*{2. Definite implementation mismatch: the economic simulation is not incident-level paired}

\paragraph{Location.}
\emph{Main\_economic.tex}, ``Incident-level paired counterfactual,'' especially
the sentence ``Both costs are calculated for the same simulated incident
severity, payer, condition, and underlying risk,'' and
\path{pec_economic_simulation.py}.

\paragraph{Problem.}
The executable model does not simulate incident-level
\(C_{ip}^{(0)}\) and \(C_{ip}^{(1)}\), severity, payer, or condition.  It samples
one pathway-level PERT distribution for an already-net per-changed-incident
delta and multiplies it by an expected affected count.  Payer/system and
condition-family weights are applied later as fixed reporting allocations.
The model therefore implements an aggregate pathway scenario, not the paired
micro-simulation described in the methods.

The same structure also leaves cross-pathway overlap unresolved for economic
outcomes.  The acute pathways are valued separately against overlapping
opportunity populations, and the scalar \(w_j\) attenuation does not allocate
an actual incident to a unique economic pathway.  Feature weights merely
redistribute value after pathway totals have already been added.

\paragraph{Why it matters.}
The stated pairing is offered as protection against mismatched severity and
double counting.  Because that pairing is absent, the methods overstate the
model's granularity, and the uncertainty analysis does not include the claimed
incident-level variation in payer, condition, and severity.

\paragraph{Suggested fix.}
Choose one of two defensible approaches.  Either (i) revise the paper throughout
to call the executable model an aggregate pathway-level expected-delta model
and remove claims about simulated incident pairing, or (ii) implement a true
paired incident model with shared incident attributes, explicit baseline/PEC
resource ledgers, and one cross-pathway allocation or joint-outcome mechanism.
In either case, define exactly what \(w_j\) removes and show that the remaining
pathway totals are non-overlapping.

\subsection*{3. Definite algorithmic mismatch: acute annual counts are not mutually exclusive across pathways}

\paragraph{Location.}
\emph{main.tex}, ``Mature full-implementation model'' and ``Mature
probabilistic results,'' especially Equations~\((16)\) and \((18)\) in the
manuscript; \path{pec_full_potential_simulation.py}, the
\texttt{outcome\_noise} branch.

\paragraph{Problem.}
The code normalizes the five acute relevance shares so that their sum does not
exceed one, which is correct for expected values.  It then draws a separate
binomial affected count for each pathway from the same
\(36.367261\)-million incident denominator.  Independent binomial draws do not
make the pathway memberships mutually exclusive: the same conceptual incident
can appear in more than one pathway's count realization.  The within-pathway
benefit/harm/none draw is multinomial-equivalent, but the cross-pathway category
assignment is not.

\paragraph{Why it matters.}
This contradicts the statements that ``one mortality outcome is allocated once
across all five dominant mechanisms'' and that the five mechanisms are
mutually exclusive.  The numerical effect is currently small because parameter
uncertainty dominates annual count noise, but the predictive variance and
cross-pathway covariance are not generated by the claimed model.

\paragraph{Suggested fix.}
For each parameter draw, first draw one multinomial allocation across the five
acute categories plus an unaffected remainder.  Apply pathway-specific success
and benefit/harm/none outcomes within those allocated counts.  Use a separate
random-number stream for annual count noise so that enabling or disabling count
noise cannot alter the underlying parameter draws.

\subsection*{4. Definite implementation mismatch: the printed fire-loss equation is not executable}

\paragraph{Location.}
\emph{Main\_economic.tex}, ``Location and responder access,'' Equation~\((9)\)
and the statement that the property result is subject to its incident-level
cap; Appendix ``Pathway economic deltas''; and
\path{pec_economic_simulation.py}.

\paragraph{Problem.}
The manuscript defines
\[
B_{\mathrm{fire}}=N_fL_f\left[1-\exp(-\eta\Delta t)\right],
\]
but the code contains no corresponding fire count \(N_f\), baseline per-fire
loss \(L_f\), time recovery \(\Delta t\), damage-growth parameter \(\eta\),
exponential calculation, or cap against the national fire-loss base.  Instead,
it samples a generic property delta for every affected location/access event
(and separate deltas for a few other pathways) and multiplies by a common
property driver.

\paragraph{Why it matters.}
The \(\$108\)-million property headline is not generated by the causal/time
formula used to justify it.  Readers cannot determine how many modeled events
are fires, how the cited national loss totals constrain the output, or whether
the same property benefit is being applied to nonfire location events.

\paragraph{Suggested fix.}
Either implement Equation~\((9)\) with an explicit fire-event subset and an
executable national-loss cap, or delete the equation and describe the actual
per-event PERT model.  Report the fire-event fraction, time-saved distribution,
elasticity prior, and cap diagnostic in the machine-readable output.

\subsection*{5. Definite reproducibility defect: seeds and companion mortality streams are inconsistent}

\paragraph{Location.}
\emph{Main\_economic.tex}, ``Monte Carlo simulation,'' ``Mortality-risk
valuation,'' and ``Reproducible files''; \path{pec_economic_simulation.py};
\path{pec_full_potential_simulation.py}.

\paragraph{Problem.}
The economic paper says that each scenario uses random seed \(20260808\).  The
script actually calls the current scenario with offset 101 and the mature
scenario with offset 202, producing economic seeds \(20260909\) and
\(20261010\).  The JSON records only the base seed, not these actual scenario
seeds.

The economic model also regenerates mature domestic mortality with
\texttt{outcome\_noise=False}, whereas the mortality paper's primary run uses
\texttt{outcome\_noise=True}.  Because annual count draws are interleaved with
later pathway parameter draws in one random-number stream, the parameter arrays
change when outcome noise is disabled.  Consequently, the economic paper's
domestic mature parameter median is \(3536.826\) lives (reported as
\(3{,}537\)), while the mortality paper's parameter median is \(3534.913\)
(reported as \(3{,}535\)).  This is not the identical companion stream claimed
in the reproducibility discussion.

\paragraph{Why it matters.}
Exact seeds and identical imported draws are load-bearing reproducibility
claims.  The discrepancy is small numerically but shows that the random-number
architecture couples parameter uncertainty to the optional predictive-noise
calculation.

\paragraph{Suggested fix.}
Record the actual current and mature economic seeds in JSON and manuscript
text.  Split parameter, pathway, and annual-count random streams using named
\texttt{SeedSequence} children.  Prefer importing archived companion path
arrays or a parameter-only stream whose output is invariant to whether count
noise is requested.  Reconcile all \(3{,}535\)/\(3{,}537\) and related Tier-C
figures after the fix.

\subsection*{6. Definite numerical inconsistency: the ten-year table is stale}

\paragraph{Location.}
\emph{Main\_economic.tex}, the ten-year projection table, the
executive-summary row for ten-year present value, and
\path{pec_economic_results.json}.

The executive summary correctly reports the mature VSL-inclusive \(7\%\)
present value as \(\$237.132\) billion, but the ten-year table reports
\(\$237.063\) billion.  The rerun gives the following societal totals:

\begin{center}
\centering
\small
\begin{tabular}{@{}lrr@{}}
\toprule
Measure & Excluding VSL & Including VSL\\
\midrule
Cumulative constant 2026 dollars & \(\$69.993\)b & \(\$374.916\)b\\
Cumulative nominal dollars & \(\$81.223\)b & \(\$435.072\)b\\
Present value, \(1.6\%\) real & \(\$62.716\)b & \(\$335.936\)b\\
Present value, \(3\%\) real & \(\$57.116\)b & \(\$305.941\)b\\
Present value, \(7\%\) real & \(\$44.270\)b & \(\$237.132\)b\\
\bottomrule
\end{tabular}

\emph{Values reproduced from the current executable output; these should
replace the stale societal columns in the manuscript table.}
\end{center}

\paragraph{Why it matters.}
The discrepancy is small but directly contradicts the claimed reproducible
link between the paper and output JSON.

\paragraph{Suggested fix.}
Generate the ten-year table from JSON rather than copying rounded values by
hand, and add a regression assertion for every printed projection total.

\subsection*{7. Definite archive defect: the claimed SHA-256 manifest is absent}

\paragraph{Location.}
Both manuscripts' reproducibility sections and
\path{source-materials/README.md} refer to
\path{reproducibility-manifest.sha256} at project root.

\paragraph{Problem.}
That file is not present in the supplied archive.  The manuscripts therefore
make a false archive-completeness claim even though the main scripts, JSON, and
CSV outputs are present and rerunnable.

\paragraph{Why it matters.}
The missing manifest prevents the promised verification that the manuscripts,
inputs, scripts, outputs, and source transcriptions form one exact archived
release.

\paragraph{Suggested fix.}
Generate the manifest after all manuscript and result corrections, include it
in the archive, and add a documented verification command.  The manifest
itself must not be regenerated silently without updating the release date or
version.

\subsection*{8. Likely issue: the CPR odds-ratio conversion uses an incompatible marginal baseline}

\paragraph{Location.}
\emph{main.tex}, ``Clinically eligible OHCA time-response conversion,''
Equations~\((3)\)--\((4)\) and the statement that \(15{,}000/78{,}048=19.2\%\)
is the ``compatible baseline.''

\paragraph{Problem.}
The cited study's \(19.2\%\) is overall survival across all bystander-CPR delay
categories.  The reported reference-category survival for CPR within one
minute is \(22.4\%\), and the \(0.91\) and \(0.73\) values are adjusted
category odds ratios, not estimates from a fitted constant causal log-odds
slope.  Dividing \(-\log(0.91)\) by two and \(-\log(0.73)\) by four and applying
that slope to the marginal \(19.2\%\) risk is therefore a planning
approximation, not an exact conversion from a compatible baseline.

\paragraph{Why it matters.}
This affects the implementation-facing interpretation of
\(\beta_{\mathrm{CPR}}\).  The current language component is small, but the same
method could become material if the eligible population or recovered-time
assumptions increase.

\paragraph{Suggested fix.}
Use category-specific risks and uncertainty, fit a curve to the reported delay
categories, or obtain standardized risks from the authors.  At minimum, replace
``compatible baseline'' with ``marginal same-cohort baseline used for an
approximate transport,'' and state that adjusted odds ratios are being treated
as a constant per-minute planning slope.

\subsection*{9. Role/terminology drift: Functions 1--5 have conflicting mortality descriptions}

\paragraph{Location.}
\emph{main.tex}, ``Treatment of all 23 functions'' and the mature ``Earlier
correct 911'' pathway; \emph{Main\_economic.tex}, Table~1.

\paragraph{Problem.}
The mortality function table says Functions 1--6 contribute ``only through''
the whole-program 988 attribution.  The mature mortality model nevertheless
contains an ``Earlier correct 911'' pathway, while the economic paper explicitly
maps Functions 1--5 to correct first-service selection and earlier 911 access.
The mature mortality paper does not state which functions feed the earlier-911
row, leaving the same functions with conflicting role descriptors across the
manuscripts.

\paragraph{Why it matters.}
The ambiguity prevents an implementer from tracing all 23 functions into the
mature mortality model and weakens the claim of full integration.

\paragraph{Suggested fix.}
Provide one cross-paper function-to-pathway matrix with separate columns for
the evidence-only, current-attribution, mature mortality, and economic models.
If Functions 1--5 feed mature earlier-911 mortality, revise ``only through
988'' to apply explicitly to the current-attribution scenario.

\section*{Claim calibration and meaningful editorial findings}

\subsection*{Mortality-neutral is not evidence-neutral}

\paragraph{Location.}
\emph{Main\_economic.tex}, executive summary and conclusion, which call
\(\$8.35\) billion the ``clean mortality-neutral result.''

\paragraph{Problem and significance.}
The sensitivity removes beneficial and adverse mortality terms, but it retains
author-specified medical-resource deltas, nonfatal productivity, QALYs,
disability effects, adoption, relevance, and technical-success priors.  The
result is neutral only with respect to mortality credit, not with respect to
clinical or economic evidence.

\paragraph{Suggested fix.}
Use ``mortality-credit-neutral planning result'' and preserve the same
assumption-conditional caveat used for the other economic totals.

\subsection*{Notation overload should be removed}

\paragraph{Location.}
\emph{Main\_economic.tex} uses \(G_{ir}\) for dollar-valued gross public
resource and \(G_j\) for a dimensionless maturity multiplier.
\emph{main.tex} uses \(A_j\) for affected incidents while also using
unsubscripted \(A\) for the shared user-activation multiplier.

\paragraph{Why it matters.}
These are distinct quantities with different units and roles, and the overload
is avoidable in implementation-facing equations.

\paragraph{Suggested fix.}
Rename the economic maturity multiplier \(M_j\), and rename the shared
activation factor \(U\) or \(A_{\mathrm{act}}\).  Add units to the symbol
definitions.

\section*{Proofing sweep}

The bundled mechanical scan was run on both compiled manuscripts.  Its
semicolon-capitalization hits were false positives caused by acronyms, proper
nouns, table grade labels, and publication titles.  It found no verified
duplicated punctuation, malformed ``into \ldots{} into'' frames, product-name
capitalization errors, or \(\arctan(x/y)\)-type branch ambiguity.  No inverse
trigonometric expressions occur in the reviewed equations.

The following high-confidence cleanup items remain:
\begin{itemize}
\item \textbf{Bibliographies:} several entries are incomplete.  These include
\texttt{kreuter211}, \texttt{ehr2024}, \texttt{fall2022},
\texttt{fallstudy}, \texttt{autocall2023}, \texttt{autocallstudy},
\texttt{hia2023}, and \texttt{pif2013}.  Add authors, year, journal/report
title, volume/pages, and DOI or stable report identifier where available.
\item \textbf{Economic pathway table heading:} ``Excl. VSL/ mortality
production'' is ambiguous.  Replace it with ``Excluding VSL and mortality
production'' to match the caption and code.
\item \textbf{Executable documentation:}
\path{pec_economic_simulation.py} says its assumptions are documented in
\path{pec_economic_value.tex}, which is not in the archive.  Replace that name
with \path{Main_economic.tex} or rename the manuscript consistently.
\end{itemize}

\section*{External-verification status}

Spot checks support the quoted primary-source values for
\texttt{nguyen2024}, \texttt{video2025}, \texttt{hhsvalues2026},
\texttt{patel2026}, and \texttt{ombm2523}.  The paper correctly labels the
988 result as an external association rather than a PEC effect.

The following items still require external or new empirical verification and
must not be upgraded from planning assumptions in the current text:
\begin{itemize}
\item the scene-to-patient mortality transport factor \(\kappa_A\);
\item a national fire response-time/property-loss elasticity matching the
proposed endpoint and implementation;
\item PEC-specific reach, activation, end-to-end success, overlap, beneficial
ARR, adverse ARR, QALY, and per-event economic deltas;
\item traveler emergency-number unfamiliarity and incident consequences by
destination; and
\item the provenance of the unavailable original PEC binary source documents.
\end{itemize}
Time-sensitive national anchors should also be refreshed as one dated release
and accompanied by archived source extracts or a data dictionary; this review
did not independently revalidate every anchor in the two long national-input
tables.

\section*{Highest-priority fixes}

\begin{enumerate}
\item Rebuild the public fiscal/capacity calculation from a common gross
resource ledger; do not rely on the current tautological assertion.
\item Decide whether the economic analysis is an aggregate pathway model or a
true paired incident simulation, and make the code and methods agree.
\item Implement one multinomial allocation for the five acute mortality
pathways and isolate parameter RNG streams from annual count noise.
\item Implement or remove the printed fire-loss equation and cap.
\item Correct seeds, companion-stream handling, the ten-year table, and the
missing SHA-256 manifest; then regenerate all manuscript tables directly from
machine-readable results.
\item Clarify the CPR odds-ratio transport and publish a single cross-paper
function-to-pathway map.
\end{enumerate}

After these changes, the manuscripts' careful identification caveats would be
substantially better matched by the executable model.  Until then, the
mortality results remain transparent planning scenarios, while the economic
headlines should be treated as provisional because the claimed non-overlap,
pairing, and property-loss mechanisms are not yet implemented as described.

\end{document}